\section{模型检验}

本文的检验分为三个层次：对单源几何结论的独立数值核验、对多源策略的配对消融与压力测试，以及对模型适用边界的显式声明。三者的作用不同——前者检查"算得对不对"，中者检查"改得有没有用"，后者界定"结论能用在哪里"。

\subsection{单源几何的独立核验}

问题一与问题二的核心是几何计算，因此本文使用\textbf{与主实现不同的方法}重新求解同一批对象，以避免共享同一处代码缺陷。

\begin{table}[H]
  \centering
  \caption{问题一、问题二的独立数值核验}
  \label{tab:verify-q12}
  \begingroup
  \renewcommand{\arraystretch}{1.4}
  \begin{tabularx}{\textwidth}{L{0.24\textwidth}C{0.14\textwidth}YC{0.19\textwidth}}
    \toprule[1.2pt]
    {\heifont\bfseries 项目} & {\heifont\bfseries 检查项数} & {\heifont\bfseries 独立方法} & {\heifont\bfseries 最大差异} \\
    \midrule[0.5pt]
    问题一定位区域 & 260 项 & 线性规划支持函数核对枚举顶点 & $4.55\times10^{-13}$ 米 \\
    问题一最小包围圆 & 含于上项 & 缩放凸范数优化核对支撑圆 & $5.96\times10^{-10}$ 米 \\
    问题二后验区域 & 1305 项 & 另一构造 + 线性规划支持函数 & 见对应运行记录 \\
    问题二方向工况 & 1100 条 & 11 种距离 $\times$ 5 种首次误差 $\times$ 5 种第二次误差 & 全部未失信号 \\
    \bottomrule[1.2pt]
  \end{tabularx}
  \endgroup
\end{table}

问题一的核验包含 3500 次直接 $\arctan2$ 角度判断、60 组多站随机一致观测、解析特例，以及空集、无界、点、线段四类退化情形与平移旋转不变性检查。问题二另有 300 组、约 120 万个直接位置检查支持其实现。

需强调这些差异是该验证集上的数值差异，\textbf{不是}对任意输入的误差上界；浮点求交也不是区间算术意义下的严格证明。

\subsection{多源策略的配对审计与消融}

第三、四问的比较一律采用同场配对：同一场景分别运行对照版本与本模型，比较逐场的 $T/N$。这样做的目的是避免把场景差异误认为策略差异。

第三问的消融结果如表 \ref{tab:ablation-q3} 所示。

\begin{table}[H]
  \centering
  \caption{第三问消融实验（20 场，场景等权 $T/N$，单位：秒/点）}
  \label{tab:ablation-q3}
  \begingroup
  \renewcommand{\arraystretch}{1.4}
  \begin{tabularx}{\textwidth}{L{0.38\textwidth}C{0.18\textwidth}C{0.18\textwidth}Y}
    \toprule[1.2pt]
    {\heifont\bfseries 配置} & {\heifont\bfseries 时间/秒} & {\heifont\bfseries 差值/秒} & {\heifont\bfseries 结论} \\
    \midrule[0.5pt]
    不更新失败光学排除区域 & 238.0196 & -- & 对照 \\
    更新失败光学排除区域   & 237.9988 & $-0.0208$ & 区间 $[-0.3224,0.3640]$ 含 0 \\
    \bottomrule[1.2pt]
  \end{tabularx}
  \endgroup
\end{table}

该结果表明：失败信息更新在\textbf{几何上是有效的}（它确实删除了已被证书排除的区域），但其对整场效率的贡献\textbf{尚未被证实}，不能把 v6 的整体改善归因于这一单项。

第四问的消融以"覆盖布局"为核心变量。40 场配对显示布局替换带来 26.00\% 的平均改善，但其中路线与可见性选点同时变化，因此\textbf{当前没有足够的消融证据把全部收益归于单一模块}。48 场压力测试（边界外向、边界切向、近源聚集、不同定向源比例）全部完整清除且各组均值改善，但压力组是构造样例，不能与普通随机场景混合后宣称平均达标。

\subsection{稳健性与敏感性}

\textbf{误差界敏感性。} 由于题面未完全明确 $\pm1^\circ$ 是对舍入前还是最终数值的保证，本文另设 $\varepsilon=1.005^\circ$ 的保守口径（覆盖"先有 $\pm1^\circ$ 误差、再四舍五入到 0.01$^\circ$"的情形）。第三、四问的实现均采用该保守口径，因此结论不依赖这一歧义的取舍。

\textbf{离散化敏感性。} 圆盘一律用外接正多边形代替，本文用 180、360、720 三种边数重复检查，确认定位半径评价对离散分辨率不敏感。需注意单圆的外包偏差不是交集误差的上界，尤其在几何病态时不能直接传播。

\textbf{误差场压力。} 在固定极端测角误差（恒定 $+1^\circ$、恒定 $-1^\circ$）与棋盘式误差场下，第三问的 v6 优势消失（分别约慢 0.48、0.30 秒/点，棋盘组基本持平），第四问则仍有改善。这说明效率提升依赖误差的实际空间结构，不能推广到未知的官方误差场。

\textbf{场景规模的稳健性。} 按源数分层的结果显示，源数越少每源摊销的搜索成本越高（第三问 10 源场景约为 16 源场景的 1.45 倍），因此不能用一个混合均值掩盖小源数场景的劣势，本文保留分层结果。

\subsection{模型适用边界的显式声明}

本文认为，明确写出模型\textbf{不能}用于何处，与给出结果同样重要。以下边界在正文各处已经出现，此处集中列出。

\begin{itemize}[label=\textbullet,leftmargin=2em]
  \item \textbf{保证的层次。} 问题一的 $D\le40$ 不是保证清除的充分条件，正确条件是 $\varrho(P)\le20$；问题二的三圆盘区域是通用\textbf{充分}区域，未声称是其考虑全部约束后的最大可行域；问题二的推荐点是有限候选搜索的结果，既没有证明连续空间全局最优，也没有对未入围粗点的最终精度分值全部求解。
  \item \textbf{全向前提不可移植。} 问题二的三圆盘接收保证与问题三的无信号圆盘排除、全向覆盖证书，\textbf{都依赖全向接收}，不能不加修改地用于第四问的定向源。
  \item \textbf{信息使用的界限。} 所有执行策略只使用已实际执行的观测；隐藏真值仅用于离线生成响应与独立验证。到达某处不等于测过全部频道；计划扫描不等于已完成覆盖。
  \item \textbf{终止条件的界限。} 源数未知，已发现 10 个或 14 个源都不足以结束；已发现 16 个频道可以停止继续搜索未知频道，但仍须清除全部已发现源。
  \item \textbf{数值推广的界限。} 官方演练与本地合成实验分开记录；本地计费与官方日志匹配，\textbf{不能}推出布局分布、接收半径分布、定向源比例与固定误差场也匹配。压力案例是构造样例，不能与普通随机场景混合后声称平均达标。
  \item \textbf{最优性的界限。} 连续调整扫描点的子问题虽然可解，但整体组合问题非凸，路径重排、删点与移点只是交替改进，不构成全局最优证明；圆周 10 源的 233.64 秒/点只否定"所有合法场景上的 200 秒硬保证"，不否定特定分布下均值可能更低。
\end{itemize}

\subsection{尚待完成的检验}

按题目要求，第三、四问需在充分演练后进行三次正式测试，并按表 1 格式报告结果、导出日志文件。\textbf{这两组正式测试目前均未进行}，因此本文第三、四问的全部数值均属本地合成实验与一次官方演练的记录，不构成正式成绩。测试完成后，应补充的内容包括：三次测试的案例编码、清除干扰源个数、平均定位清除时间与程序运行时间，以及对应的模拟器日志文件。这部分内容留待测试结束后填写。
